\section{関連研究}\label{sec:related}

\subsection{頻出パターンマイニング}

Aprioriアルゴリズム~\cite{agrawal1994fast}は、
反単調性（anti-monotonicity）を利用した候補生成・剪定戦略により、
頻出アイテムセットを効率的に列挙する手法である。
ECLAT~\cite{zaki2000scalable}は垂直データ表現を用い、
集合の共通部分演算によってサポートを計算する。
これらの手法はトランザクション全体にわたる頻出パターンを検出するが、
時間的な局所性（temporal locality）は考慮しない。

密集区間マイニングは、スライディングウィンドウ内でのサポート集中を検出することで、
時間的局所性を捉える。
しかし、検出された密集区間の将来予測には踏み込んでいない。

\subsection{バースト検出}

Kleinberg~\cite{kleinberg2003bursty}は、
テキストストリームにおけるバースト構造を隠れマルコフモデルで検出する手法を提案した。
この手法はバーストの階層的構造を発見できるが、
将来のバースト発生を予測するモデルは含まない。

\subsection{点過程}

Hawkes過程~\cite{hawkes1971spectra}は、
過去のイベントが将来のイベント発生確率を高める自己励起点過程である。
条件付き強度関数は次式で定義される：
\begin{equation}
  \lambda(t) = \mu + \sum_{t_i < t} \alpha \beta \exp(-\beta(t - t_i))
\end{equation}
ここで$\mu$はベースライン強度、$\alpha$は自己励起パラメータ、$\beta$は減衰率である。
Hawkes過程はソーシャルネットワーク分析~\cite{zhou2013learning}や
金融市場の微細構造分析に広く応用されている。

Ogata~\cite{ogata1981lewis}のthinningアルゴリズムは、
Hawkes過程からのシミュレーション（サンプリング）手法を提供する。
Du ら~\cite{du2016recurrent}は、リカレントニューラルネットワークを用いた
マーク付き時間的点過程を提案し、イベント系列の埋め込み表現を学習している。

本研究では、密集区間の発生をHawkes過程で表現し、
パターンマイニングと点過程の接点を確立する。

\subsection{生存分析}

Cox比例ハザードモデル~\cite{cox1972regression}は、
イベント発生までの時間を共変量で説明するセミパラメトリック手法である。
Kaplan-Meier推定量~\cite{kaplan1958nonparametric}は、
非パラメトリックな生存関数の推定を提供する。
Weibull分布~\cite{weibull1951statistical}は、
故障時間分析に広く用いられるパラメトリック生存モデルである。

本研究では、密集区間の持続時間をWeibull分布でモデル化し、
将来の密集区間がどの程度持続するかを予測する。

\subsection{時系列予測}

Prophet~\cite{taylor2018forecasting}やN-BEATS~\cite{oreshkin2020nbeats}は、
連続値の時系列予測に高い性能を示す。
しかし、これらの手法は連続的な値の予測を目的としており、
離散的な「区間の発生」を直接扱うことは困難である。
本研究はこれらの手法と相補的な位置づけにある。
